\documentclass[a4paper]{article}
\usepackage[fontset=fandol]{ctex}
\usepackage[margin=1.65cm]{geometry}
\usepackage{amsmath,amssymb}
\usepackage{graphicx}
\usepackage{hyperref}
\usepackage{longtable,booktabs}
\usepackage{xcolor}
\usepackage{enumitem}
\usepackage{listings}
\lstset{basicstyle=\ttfamily\small,breaklines=true,columns=fullflexible,keepspaces=true}
\hypersetup{colorlinks=true,linkcolor=blue,urlcolor=blue}
\setlist[itemize]{leftmargin=1.25em,itemsep=0.18em}
\setlength{\parskip}{0.45em}
\setlength{\parindent}{2em}
\sloppy
\title{第 18 讲：Guest Lecture: Dan Fu 与推理引擎研究}
\author{CS336 Spring 2026 public videos and official course materials}
\date{}
\begin{document}
\maketitle
\tableofcontents
\newpage
\section{本章导读}
本章讨论 \textbf{Guest Lecture: Dan Fu 与推理引擎研究}。CS336 的第一层目标不是把现成大模型当作黑箱使用，而是把 language model 从输入文本、训练目标、模型结构、数据、系统和评估这些部件重新拆开。只有拆开之后，读者才能判断一个设计选择究竟改变了什么。
设想训练好的模型要接入真实流量。Serving system（服务系统）需要调度请求、管理 KV cache、写高效 kernels、处理 prefill/decode，并在 latency、throughput 和成本之间做工程取舍。
本章对应 CS336 Spring 2026 第 18 个公开视频，并结合课程页提供的可解析资料。视频和课程页给出事实边界；正文把这些材料改写为可自学的教材章节。
阅读本章时，先理解问题为什么存在，再看定义、公式和实现形态，随后通过例题和练习检查自己是否能独立复现这条 reasoning chain。
\begin{figure}[h]
\centering
\includegraphics[width=0.58\linewidth]{source_anchor.jpg}
\caption{第 18 讲的课程图像锚点。}
\end{figure}
\subsection{本章知识路线}
\begin{itemize}
\item 本章首先说明 Guest Lecture: Dan Fu 与推理引擎研究 为什么是 language modeling pipeline 中不可跳过的一环。
\item 其次定义本章反复使用的术语，并说明这些术语对应的计算对象或系统对象。
\item 随后用公式和伪代码表达主要机制，让概念能够落到可计算的变量上。
\item 最后给出例题、风险讨论和练习，便于读者检查自己是否真正掌握了机制。
\end{itemize}
\subsection{结构图}
下面的图用于帮助读者定位本章的核心结构。先看概念之间的依赖，再回到正文中的公式、伪代码和例题。
\begin{figure}[h]
\centering
\includegraphics[width=0.72\linewidth]{figures/source_grounded_concept_map.png}
\caption{第 18 讲概念图：本章问题、术语和机制的关系。}
\end{figure}
\begin{figure}[h]
\centering
\includegraphics[width=0.72\linewidth]{figures/source_grounded_checklist.png}
\caption{第 18 讲学习路径：从问题定义进入公式、实现、实验和练习。}
\end{figure}
\section{本章主线}
本章可以沿着三条线索阅读：课程为什么把这个主题放在这里，它改变了哪些变量，以及怎样用小实验检查这些变量。
\subsection{本章要解决的核心问题}
Dan Fu guest lecture 从 serving 侧看语言模型：训练好的模型要变成真实服务，需要 inference engine、GPU kernels、scheduler、KV cache 管理和生产流量处理。
这一点对应的学习动作是：找出对象，写出变量，说明变量如何被测量。若输入规模、硬件、数据分布或评估协议改变，结论也必须重新检查。
因此，Guest Lecture: Dan Fu 与推理引擎研究 不应被读成若干名词的集合，而应被读成一组可检验的机制。
\subsection{从机制到资源账本}
推理系统是 full-stack innovation 的位置。算法结构、kernel library、memory bandwidth、batching 和 request distribution 会共同决定用户体验。
这一点对应的学习动作是：找出对象，写出变量，说明变量如何被测量。若输入规模、硬件、数据分布或评估协议改变，结论也必须重新检查。
因此，Guest Lecture: Dan Fu 与推理引擎研究 不应被读成若干名词的集合，而应被读成一组可检验的机制。
\subsection{从课堂讲解到可复现实验}
ThunderKittens、megakernel、SSM/recurrence 等讨论说明 serving 不是模型训练后的附属环节，而是会反过来影响模型架构研究的问题。
这一点对应的学习动作是：找出对象，写出变量，说明变量如何被测量。若输入规模、硬件、数据分布或评估协议改变，结论也必须重新检查。
因此，Guest Lecture: Dan Fu 与推理引擎研究 不应被读成若干名词的集合，而应被读成一组可检验的机制。
\section{术语、符号与问题边界}
本章的术语横跨模型、数据和系统。读者需要先知道每个词指向的对象：有些词描述数学目标，有些词描述张量形状，有些词描述硬件瓶颈，还有一些词描述评估协议。术语边界不清，后面的公式和实现就会变成空洞的缩写。
\begin{longtable}{p{0.08\linewidth}p{0.42\linewidth}p{0.42\linewidth}}
\toprule
\textbf{\#} & \textbf{术语} & \textbf{阅读要求}\\
\midrule
1 & serving system（服务系统） & 给出定义、一个最小例子、一个关联公式或代码位置，以及一个典型失败情形。\\
2 & inference engine（推理引擎） & 给出定义、一个最小例子、一个关联公式或代码位置，以及一个典型失败情形。\\
3 & GPU kernel & 把它放入资源账本：单位是什么，随输入规模怎样变化，最容易被哪类 benchmark 误读。\\
4 & ThunderKittens & 给出定义、一个最小例子、一个关联公式或代码位置，以及一个典型失败情形。\\
5 & decode loop（解码循环） & 给出定义、一个最小例子、一个关联公式或代码位置，以及一个典型失败情形。\\
6 & state-space model (SSM) & 给出定义、一个最小例子、一个关联公式或代码位置，以及一个典型失败情形。\\
\bottomrule
\end{longtable}
这些术语之间有依赖关系。例如 tokenization 改变 sequence length，sequence length 又改变 attention FLOPs、KV cache 和 evaluation token budget；parallelism 改变显存占用，也改变通信瓶颈和故障恢复方式。
\subsection{关键词}
下面的关键词在课程材料中反复出现。它们不代替定义，只提示读者本章哪些对象需要在后文持续跟踪。
\begin{itemize}
\item executable, linked, public, playlist, item, spring, page.
\end{itemize}
\section{从训练到 serving 的另一侧}
现在进入 从训练到 serving 的另一侧。它把 Guest Lecture: Dan Fu 与推理引擎研究 中的一部分问题推进到可操作层面：哪些对象要被定义，哪些变量要被测量，哪些限制会改变结论。
\subsection{为什么要讨论 从训练到 serving 的另一侧}
Dan Fu guest lecture 的视频开场说明：课程大多讲如何训练模型，本讲从拥有模型后的另一侧看 inference 和 serving。
推理引擎把抽象的模型计算图映射到具体 GPU kernels、调度器、cache 管理和生产流量。
字幕中反复出现 inference、GPU、kernel、decode、cache、request 等词，显示核心主题是 serving stack。
从训练到 serving 的另一侧 的作用是把 Guest Lecture: Dan Fu 与推理引擎研究 中的一个抽象问题变成可计算对象：哪些变量进入公式，哪些状态进入代码，哪些条件会让结果失效。
\subsection{从训练到 serving 的另一侧 的机制}
Dan Fu guest lecture 的视频开场说明：课程大多讲如何训练模型，本讲从拥有模型后的另一侧看 inference 和 serving。 这句话把讨论落在硬件和系统状态上。读者需要知道哪些量可直接观察，哪些量只能通过日志、profile 或评估协议间接推断。
推理引擎把抽象的模型计算图映射到具体 GPU kernels、调度器、cache 管理和生产流量。 因而同一个方法在不同规模下可能改变不同的变量，尤其是吞吐、延迟、显存峰值、通信量和 kernel 利用率。比较方法时，必须同时给出问题规模和实验条件。
字幕中反复出现 inference、GPU、kernel、decode、cache、request 等词，显示核心主题是 serving stack。 这使本节具有可检验性：一个结论若不能在公式、伪代码、profile trace、benchmark 或 ablation 中表现出来，就仍然只是直觉。
因此，从训练到 serving 的另一侧 应当被理解为可以实现和测试的机制，而不是一个章节标题。CS336 反复强调 from scratch，正是因为只有能被复现的机制，才可能被稳定改进。
\subsection{从训练到 serving 的另一侧 的数学表达}
\[\text{user latency}=t_{\text{queue}}+t_{\text{prefill}}+t_{\text{decode}}+t_{\text{postprocess}}\]
符号说明：用户感知延迟由排队、预填充、逐 token 解码和后处理组成；每一项都可能成为瓶颈。
这个关系式应被读成 从训练到 serving 的另一侧 的资源账本。它不承诺精确预测每次运行时间，而是帮助判断主要成本来自计算、内存访问、通信还是调度。
例如，validation loss 常被用作模型质量的 proxy，tokens/s 用作吞吐 proxy，Jaccard 或 MinHash 用作近重复 proxy，reward model score 用作偏好 proxy。proxy 不是目标本身；它只在评估协议清楚时才可引用。
\subsection{从训练到 serving 的另一侧 在系统中的实现}
把 从训练到 serving 的另一侧 写成程序时，最重要的是暴露状态变化和数据移动。下面的伪代码保留执行顺序，省略框架样板。
\begin{lstlisting}[language=Python]
request = receive_prompt()
kv = prefill(request.prompt)
for step in decode_loop:
    token, kv = engine.decode(kv)
    stream_to_user(token)
\end{lstlisting}
这段伪代码的读法，是找出 从训练到 serving 的另一侧 改变了哪些状态，以及这些状态如何对应到 吞吐、延迟、显存峰值、通信量和 kernel 利用率。如果某个状态没有被记录，后续实验就很难解释异常结果。
\subsection{边界条件与常见误解}
\begin{itemize}
\item 小规模能工作的 serving code 到大规模流量下会暴露罕见 bug。
\item 生产 workload 的输入/输出长度分布和教学 benchmark 差异很大。
\item 不要把小实验中的局部结论自动外推到 frontier-scale；scale、数据分布和硬件拓扑都会改变结论。
\item 不要把 benchmark 或 profile 的单个数字当作机制解释；必须同时记录实验设置、输入形状、版本和评估口径。
\end{itemize}
\subsection{与本章其他概念的关系}
从训练到 serving 的另一侧 与本章其他部分的关系是互补的：术语表给出对象名称，公式给出最小数学结构，伪代码给出执行顺序，例题则把这些内容压缩成一个可检查的计算。
对这一类主题，最常见的误解是只列方法名。更有用的做法是比较每个方法改变的成本项、依赖的假设和可能的失败模式。
\subsection{怎样检验 从训练到 serving 的另一侧}
检验 从训练到 serving 的另一侧 时，先把抽象说法落到硬件和系统状态上。与 inference、serving、kernel、decode、request 相关的对象必须能被构造、打印、计数或评分；否则实验只能得到描述性观察，不能得到可复现结论。
一个合适的小实验应当保留本节机制，同时让 吞吐、延迟、显存峰值、通信量和 kernel 利用率 中至少一个量发生可解释变化。输入规模不必逼真；关键是读者能手算或打印关键中间量，并说明哪个变量触发了结果变化。
随后把公式说明转成可记录的数量关系。符号说明：用户感知延迟由排队、预填充、逐 token 解码和后处理组成；每一项都可能成为瓶颈。 如果数量只能间接观测，就必须说明使用了什么 proxy；如果使用 benchmark 或 reward，也要说明协议和随机性。
对照实验只应改变一个主要因素，例如 batch size、sequence length、vocabulary size、attention heads、filter threshold、learning rate、decoding 参数或 verifier。多个因素同时变化时，实验最多说明系统表现不同，不能说明是哪一个机制在起作用。
最后要主动寻找失败条件。教材中的成功路径容易记住，但工程中真正决定方法边界的是失败案例。本节至少要记住这些限制：小规模能工作的 serving code 到大规模流量下会暴露罕见 bug；生产 workload 的输入/输出长度分布和教学 benchmark 差异很大。复习时可以把这些限制改写成反例测试。
实验结论还要放回 Guest Lecture: Dan Fu 与推理引擎研究 的整体结构中。从训练到 serving 的另一侧 会影响训练目标、数据选择、系统成本或评估解释中的至少一项；能说清楚这种影响，才说明读者已经从名词记忆进入机制理解。
\subsection{常见混淆：从训练到 serving 的另一侧}
本节容易混淆的关键词包括 inference、serving、kernel、decode、request。它们在课程材料中可能连续出现，但层级并不相同：有的命名对象，有的描述操作，有的只是指标或约束。
其中，inference（推理）是已训练模型服务请求并生成 token 的过程；serving system（服务系统）把模型计算图变成可承载真实流量的推理引擎。
区分这些词时，可以先问它们是否改变 吞吐、延迟、显存峰值、通信量和 kernel 利用率。如果一个词不能对应到变量、状态、代码路径或评估协议，它在本节中就还没有形成可操作含义。
因此，复习 从训练到 serving 的另一侧 时不应只抄关键词，而应写出至少一个最小例子：输入是什么，哪一步使用这个词，输出或指标怎样变化，失败条件在哪里出现。
\subsection{从训练到 serving 的另一侧 的小结}
\begin{itemize}
\item 讨论 从训练到 serving 的另一侧 时，应同时报告问题规模、实验设置和主要观测变量，尤其是 吞吐、延迟、显存峰值、通信量和 kernel 利用率。
\item 如果一个方法声称降低主要成本，要追问成本是否转移到了显存、通信、数据质量、训练稳定性或评估复杂度。
\item 公式和伪代码提供推理骨架；结论仍需要 profiler、ablation 或 held-out evaluation 支持。
\end{itemize}
\section{GPU kernels 与 full-stack innovation}
现在进入 GPU kernels 与 full-stack innovation。它把 Guest Lecture: Dan Fu 与推理引擎研究 中的一部分问题推进到可操作层面：哪些对象要被定义，哪些变量要被测量，哪些限制会改变结论。
\subsection{为什么要讨论 GPU kernels 与 full-stack innovation}
讲者强调理解 inference engines 和 GPU kernels 可以开启 full-stack innovation。
ThunderKittens 被作为更低层 kernel-writing library 的例子，用于获得更细粒度控制和高带宽利用。
Megakernel、fusion、cache-aware scheduling 等方向说明算法和系统可以共同设计。
系统部分的核心问题是把数学计算映射到硬件。一个算法的复杂度公式只能给出方向；真正的速度还取决于内存层级、通信拓扑、kernel shape 和调度方式。
\subsection{GPU kernels 与 full-stack innovation 的机制}
讲者强调理解 inference engines 和 GPU kernels 可以开启 full-stack innovation。 这句话把讨论落在硬件和系统状态上。读者需要知道哪些量可直接观察，哪些量只能通过日志、profile 或评估协议间接推断。
ThunderKittens 被作为更低层 kernel-writing library 的例子，用于获得更细粒度控制和高带宽利用。 因而同一个方法在不同规模下可能改变不同的变量，尤其是吞吐、延迟、显存峰值、通信量和 kernel 利用率。比较方法时，必须同时给出问题规模和实验条件。
Megakernel、fusion、cache-aware scheduling 等方向说明算法和系统可以共同设计。 这使本节具有可检验性：一个结论若不能在公式、伪代码、profile trace、benchmark 或 ablation 中表现出来，就仍然只是直觉。
因此，GPU kernels 与 full-stack innovation 应当被理解为可以实现和测试的机制，而不是一个章节标题。CS336 反复强调 from scratch，正是因为只有能被复现的机制，才可能被稳定改进。
\subsection{GPU kernels 与 full-stack innovation 的资源账本}
\[\text{bandwidth utilization}=\frac{\text{bytes served per second}}{\text{peak memory bandwidth}}\]
符号说明：decode 阶段常 memory-bandwidth bound，因此带宽利用率直接影响 tokens/s。
这个关系式应被读成 GPU kernels 与 full-stack innovation 的资源账本。它不承诺精确预测每次运行时间，而是帮助判断主要成本来自计算、内存访问、通信还是调度。
例如，validation loss 常被用作模型质量的 proxy，tokens/s 用作吞吐 proxy，Jaccard 或 MinHash 用作近重复 proxy，reward model score 用作偏好 proxy。proxy 不是目标本身；它只在评估协议清楚时才可引用。
\subsection{GPU kernels 与 full-stack innovation 在系统中的实现}
把 GPU kernels 与 full-stack innovation 写成程序时，最重要的是暴露状态变化和数据移动。下面的伪代码保留执行顺序，省略框架样板。
\begin{lstlisting}[language=Python]
load_model_tiles()
load_kv_cache_tiles()
run_fused_decode_kernel()
write_next_token_logits()
\end{lstlisting}
这段伪代码的读法，是找出 GPU kernels 与 full-stack innovation 改变了哪些状态，以及这些状态如何对应到 吞吐、延迟、显存峰值、通信量和 kernel 利用率。如果某个状态没有被记录，后续实验就很难解释异常结果。
\subsection{边界条件与常见误解}
\begin{itemize}
\item 更低层抽象换来性能，也带来更高 correctness/debug 成本。
\item kernel 微小错误在大规模 serving 中可能以极低概率触发，但影响严重。
\item 不要把小实验中的局部结论自动外推到 frontier-scale；scale、数据分布和硬件拓扑都会改变结论。
\item 不要把 benchmark 或 profile 的单个数字当作机制解释；必须同时记录实验设置、输入形状、版本和评估口径。
\end{itemize}
\subsection{与本章其他概念的关系}
GPU kernels 与 full-stack innovation 与本章其他部分的关系是互补的：术语表给出对象名称，公式给出最小数学结构，伪代码给出执行顺序，例题则把这些内容压缩成一个可检查的计算。
对这一类主题，最常见的误解是只列方法名。更有用的做法是比较每个方法改变的成本项、依赖的假设和可能的失败模式。
\subsection{怎样检验 GPU kernels 与 full-stack innovation}
检验 GPU kernels 与 full-stack innovation 时，先把抽象说法落到硬件和系统状态上。与 ThunderKittens、megakernel、bandwidth、fusion、GPU 相关的对象必须能被构造、打印、计数或评分；否则实验只能得到描述性观察，不能得到可复现结论。
一个合适的小实验应当保留本节机制，同时让 吞吐、延迟、显存峰值、通信量和 kernel 利用率 中至少一个量发生可解释变化。输入规模不必逼真；关键是读者能手算或打印关键中间量，并说明哪个变量触发了结果变化。
随后把公式说明转成可记录的数量关系。符号说明：decode 阶段常 memory-bandwidth bound，因此带宽利用率直接影响 tokens/s。 如果数量只能间接观测，就必须说明使用了什么 proxy；如果使用 benchmark 或 reward，也要说明协议和随机性。
对照实验只应改变一个主要因素，例如 batch size、sequence length、vocabulary size、attention heads、filter threshold、learning rate、decoding 参数或 verifier。多个因素同时变化时，实验最多说明系统表现不同，不能说明是哪一个机制在起作用。
最后要主动寻找失败条件。教材中的成功路径容易记住，但工程中真正决定方法边界的是失败案例。本节至少要记住这些限制：更低层抽象换来性能，也带来更高 correctness/debug 成本；kernel 微小错误在大规模 serving 中可能以极低概率触发，但影响严重。复习时可以把这些限制改写成反例测试。
实验结论还要放回 Guest Lecture: Dan Fu 与推理引擎研究 的整体结构中。GPU kernels 与 full-stack innovation 会影响训练目标、数据选择、系统成本或评估解释中的至少一项；能说清楚这种影响，才说明读者已经从名词记忆进入机制理解。
\subsection{常见混淆：GPU kernels 与 full-stack innovation}
本节容易混淆的关键词包括 ThunderKittens、megakernel、bandwidth、fusion、GPU。它们在课程材料中可能连续出现，但层级并不相同：有的命名对象，有的描述操作，有的只是指标或约束。
其中，kernel（内核）是高层 tensor 操作落到 GPU 的执行单元；GPU 是 CS336 systems 部分的主要硬件对象。
区分这些词时，可以先问它们是否改变 吞吐、延迟、显存峰值、通信量和 kernel 利用率。如果一个词不能对应到变量、状态、代码路径或评估协议，它在本节中就还没有形成可操作含义。
因此，复习 GPU kernels 与 full-stack innovation 时不应只抄关键词，而应写出至少一个最小例子：输入是什么，哪一步使用这个词，输出或指标怎样变化，失败条件在哪里出现。
\subsection{GPU kernels 与 full-stack innovation 的小结}
\begin{itemize}
\item 讨论 GPU kernels 与 full-stack innovation 时，应同时报告问题规模、实验设置和主要观测变量，尤其是 吞吐、延迟、显存峰值、通信量和 kernel 利用率。
\item 如果一个方法声称降低主要成本，要追问成本是否转移到了显存、通信、数据质量、训练稳定性或评估复杂度。
\item 公式和伪代码提供推理骨架；结论仍需要 profiler、ablation 或 held-out evaluation 支持。
\end{itemize}
\section{SSM/recurrence 与 inference 研究问题}
现在进入 SSM/recurrence 与 inference 研究问题。它把 Guest Lecture: Dan Fu 与推理引擎研究 中的一部分问题推进到可操作层面：哪些对象要被定义，哪些变量要被测量，哪些限制会改变结论。
\subsection{为什么要讨论 SSM/recurrence 与 inference 研究问题}
字幕中出现 state space、SSM、recurrence 和 scaling laws，说明 guest lecture 还讨论了模型结构如何影响 serving。
recurrent/SSM 风格模型有潜力改变长上下文和 decode 的状态管理方式。
但新的结构要同时满足训练稳定性、质量、硬件效率和部署复杂度。
推理阶段的关键区别在于请求是在线到达的，输出又是逐 token 生成的。训练时可以用大 batch 摊薄成本；服务时还要同时考虑首 token 延迟、吞吐、显存占用和长尾请求。
\subsection{SSM/recurrence 与 inference 研究问题 的机制}
字幕中出现 state space、SSM、recurrence 和 scaling laws，说明 guest lecture 还讨论了模型结构如何影响 serving。 这句话把讨论落在模型规模、token 数和训练预算上。读者需要知道哪些量可直接观察，哪些量只能通过日志、profile 或评估协议间接推断。
recurrent/SSM 风格模型有潜力改变长上下文和 decode 的状态管理方式。 因而同一个方法在不同规模下可能改变不同的变量，尤其是参数量、训练 token 数、compute budget、loss 外推误差和 benchmark 转化关系。比较方法时，必须同时给出问题规模和实验条件。
但新的结构要同时满足训练稳定性、质量、硬件效率和部署复杂度。 这使本节具有可检验性：一个结论若不能在公式、伪代码、profile trace、benchmark 或 ablation 中表现出来，就仍然只是直觉。
因此，SSM/recurrence 与 inference 研究问题 应当被理解为可以实现和测试的机制，而不是一个章节标题。CS336 反复强调 from scratch，正是因为只有能被复现的机制，才可能被稳定改进。
\subsection{SSM/recurrence 与 inference 研究问题 的资源账本}
\[h_t=A h_{t-1}+B x_t,\qquad y_t=C h_t\]
符号说明：这是 SSM/递推模型的简化形式；\(h_t\) 是状态，\(x_t\) 是输入，\(y_t\) 是输出。
这个关系式给出 SSM/recurrence 与 inference 研究问题 的缩放假设。它把参数量、训练 token 数和 compute budget 放到同一张账本中，重点是判断小实验能否可靠外推到大训练。
例如，validation loss 常被用作模型质量的 proxy，tokens/s 用作吞吐 proxy，Jaccard 或 MinHash 用作近重复 proxy，reward model score 用作偏好 proxy。proxy 不是目标本身；它只在评估协议清楚时才可引用。
\subsection{SSM/recurrence 与 inference 研究问题 在系统中的实现}
把 SSM/recurrence 与 inference 研究问题 写成程序后，抽象机制会变成输入、状态和输出之间的变换。下面的伪代码保留核心计算顺序。
\begin{lstlisting}[language=Python]
state = init_state()
for token in sequence:
    state = recurrent_update(state, token)
    logits = readout(state)
\end{lstlisting}
这段伪代码的读法，是找出 SSM/recurrence 与 inference 研究问题 改变了哪些状态，以及这些状态如何对应到 参数量、训练 token 数、compute budget、loss 外推误差和 benchmark 转化关系。如果某个状态没有被记录，后续实验就很难解释异常结果。
\subsection{边界条件与常见误解}
\begin{itemize}
\item 递推结构的训练并行性和长程记忆质量需要实证验证。
\item serving 系统收益只有在端到端 workload 中测量才可信。
\item 不要把小实验中的局部结论自动外推到 frontier-scale；scale、数据分布和硬件拓扑都会改变结论。
\item 不要把 benchmark 或 profile 的单个数字当作机制解释；必须同时记录实验设置、输入形状、版本和评估口径。
\end{itemize}
\subsection{与本章其他概念的关系}
SSM/recurrence 与 inference 研究问题 与本章其他部分的关系是互补的：术语表给出对象名称，公式给出最小数学结构，伪代码给出执行顺序，例题则把这些内容压缩成一个可检查的计算。
对这一类主题，最常见的误解是只列方法名。更有用的做法是比较每个方法改变的成本项、依赖的假设和可能的失败模式。
\subsection{怎样检验 SSM/recurrence 与 inference 研究问题}
检验 SSM/recurrence 与 inference 研究问题 时，先把抽象说法落到模型规模、token 数和训练预算上。与 SSM、recurrence、state、long context、decode 相关的对象必须能被构造、打印、计数或评分；否则实验只能得到描述性观察，不能得到可复现结论。
一个合适的小实验应当保留本节机制，同时让 参数量、训练 token 数、compute budget、loss 外推误差和 benchmark 转化关系 中至少一个量发生可解释变化。输入规模不必逼真；关键是读者能手算或打印关键中间量，并说明哪个变量触发了结果变化。
随后把公式说明转成可记录的数量关系。符号说明：这是 SSM/递推模型的简化形式；\textbackslash{}(h\_t\textbackslash{}) 是状态，\textbackslash{}(x\_t\textbackslash{}) 是输入，\textbackslash{}(y\_t\textbackslash{}) 是输出。 如果数量只能间接观测，就必须说明使用了什么 proxy；如果使用 benchmark 或 reward，也要说明协议和随机性。
对照实验只应改变一个主要因素，例如 batch size、sequence length、vocabulary size、attention heads、filter threshold、learning rate、decoding 参数或 verifier。多个因素同时变化时，实验最多说明系统表现不同，不能说明是哪一个机制在起作用。
最后要主动寻找失败条件。教材中的成功路径容易记住，但工程中真正决定方法边界的是失败案例。本节至少要记住这些限制：递推结构的训练并行性和长程记忆质量需要实证验证；serving 系统收益只有在端到端 workload 中测量才可信。复习时可以把这些限制改写成反例测试。
实验结论还要放回 Guest Lecture: Dan Fu 与推理引擎研究 的整体结构中。SSM/recurrence 与 inference 研究问题 会影响训练目标、数据选择、系统成本或评估解释中的至少一项；能说清楚这种影响，才说明读者已经从名词记忆进入机制理解。
\subsection{常见混淆：SSM/recurrence 与 inference 研究问题}
本节容易混淆的关键词包括 SSM、recurrence、state、long context、decode。它们在课程材料中可能连续出现，但层级并不相同：有的命名对象，有的描述操作，有的只是指标或约束。
其中，decode（解码）逐 token 生成，常受 memory bandwidth 和 KV cache 影响。
区分这些词时，可以先问它们是否改变 参数量、训练 token 数、compute budget、loss 外推误差和 benchmark 转化关系。如果一个词不能对应到变量、状态、代码路径或评估协议，它在本节中就还没有形成可操作含义。
因此，复习 SSM/recurrence 与 inference 研究问题 时不应只抄关键词，而应写出至少一个最小例子：输入是什么，哪一步使用这个词，输出或指标怎样变化，失败条件在哪里出现。
\subsection{SSM/recurrence 与 inference 研究问题 的小结}
\begin{itemize}
\item 讨论 SSM/recurrence 与 inference 研究问题 时，应同时报告问题规模、实验设置和主要观测变量，尤其是 参数量、训练 token 数、compute budget、loss 外推误差和 benchmark 转化关系。
\item 如果一个方法声称降低主要成本，要追问成本是否转移到了显存、通信、数据质量、训练稳定性或评估复杂度。
\item 公式和伪代码提供推理骨架；结论仍需要 profiler、ablation 或 held-out evaluation 支持。
\end{itemize}
\section{综合讨论：从局部机制到完整系统}
本章的主题是 Guest Lecture: Dan Fu 与推理引擎研究，但它不应被读成几个相邻知识点的拼接。从训练到 serving 的另一侧、GPU kernels 与 full-stack innovation、SSM/recurrence 与 inference 研究问题 共同回答同一个问题：语言模型系统中哪些对象可以被定义，哪些成本可以被估算，哪些结论必须通过实验或评估来确认。
这些对象包括 serving system（服务系统）、inference engine（推理引擎）、GPU kernel、ThunderKittens、decode loop（解码循环）、state-space model (SSM)。读者可以把它们看成一组接口：有的接口连接文本和 token，有的连接模型结构和张量计算，有的连接数据样本和训练目标，有的连接离线评估和在线服务。接口一旦改变，后面的成本、错误类型和优化空间也会随之改变。
\subsection{从训练到 serving 的另一侧 在系统中的位置}
从训练到 serving 的另一侧 首先提供一个局部解释：Dan Fu guest lecture 的视频开场说明：课程大多讲如何训练模型，本讲从拥有模型后的另一侧看 inference 和 serving。 这类解释的价值在于把直觉压缩成可操作对象。一旦对象明确，读者就可以追问它的输入、输出、中间状态和失败条件，而不是只记住方法名。
其次，从训练到 serving 的另一侧 会改变至少一个资源或评估变量。它可能改变 sequence length、activation size、memory bandwidth、communication volume、data mixture、reward distribution 或 benchmark protocol。这些变量在小例子里可以手算，在真实训练和服务系统里则需要 profiling、logging 和 held-out evaluation。
最后，从训练到 serving 的另一侧 不能脱离边界条件理解。一个关键 caveat 是：小规模能工作的 serving code 到大规模流量下会暴露罕见 bug。。把 caveat 写出来不是为了削弱结论，而是为了说明结论在哪些输入、硬件、数据和评估条件下才成立。
\subsection{GPU kernels 与 full-stack innovation 在系统中的位置}
GPU kernels 与 full-stack innovation 首先提供一个局部解释：讲者强调理解 inference engines 和 GPU kernels 可以开启 full-stack innovation。 这类解释的价值在于把直觉压缩成可操作对象。一旦对象明确，读者就可以追问它的输入、输出、中间状态和失败条件，而不是只记住方法名。
其次，GPU kernels 与 full-stack innovation 会改变至少一个资源或评估变量。它可能改变 sequence length、activation size、memory bandwidth、communication volume、data mixture、reward distribution 或 benchmark protocol。这些变量在小例子里可以手算，在真实训练和服务系统里则需要 profiling、logging 和 held-out evaluation。
最后，GPU kernels 与 full-stack innovation 不能脱离边界条件理解。一个关键 caveat 是：更低层抽象换来性能，也带来更高 correctness/debug 成本。。把 caveat 写出来不是为了削弱结论，而是为了说明结论在哪些输入、硬件、数据和评估条件下才成立。
\subsection{SSM/recurrence 与 inference 研究问题 在系统中的位置}
SSM/recurrence 与 inference 研究问题 首先提供一个局部解释：字幕中出现 state space、SSM、recurrence 和 scaling laws，说明 guest lecture 还讨论了模型结构如何影响 serving。 这类解释的价值在于把直觉压缩成可操作对象。一旦对象明确，读者就可以追问它的输入、输出、中间状态和失败条件，而不是只记住方法名。
其次，SSM/recurrence 与 inference 研究问题 会改变至少一个资源或评估变量。它可能改变 sequence length、activation size、memory bandwidth、communication volume、data mixture、reward distribution 或 benchmark protocol。这些变量在小例子里可以手算，在真实训练和服务系统里则需要 profiling、logging 和 held-out evaluation。
最后，SSM/recurrence 与 inference 研究问题 不能脱离边界条件理解。一个关键 caveat 是：递推结构的训练并行性和长程记忆质量需要实证验证。。把 caveat 写出来不是为了削弱结论，而是为了说明结论在哪些输入、硬件、数据和评估条件下才成立。
\subsection{把本章知识用于读论文和复现实验}
读相关论文时，可以把本章变成一个检查表。第一，论文是否清楚定义了输入规模、模型规模、数据来源和评估协议；第二，论文声称的改进落在哪个变量上；第三，作者是否报告了会让结论失效的设置；第四，实验是否能分离模型结构、数据质量、优化设置和系统实现的影响。
复现实验时，也应避免直接追求大规模。先用最小程序复现一个公式、一个伪代码分支或一个 profile 现象，再扩大输入规模。这样做的原因很简单：如果小规模程序无法解释中间量，大规模训练只会把错误隐藏在日志和随机波动里。
把这些检查合在一起，本章给出的不是一个固定答案，而是一种阅读 CS336 的方法：每个模型机制都要落到变量，每个变量都要落到测量，每个测量都要说明协议，每个协议都要承认边界。
因此，本章中的任何一句结论都可以被改写成一个可引用的完整句式：在给定数据、模型、硬件和评估协议下，某个机制改变了某个变量，并通过某个观测指标表现出来。这样的句子比单独背诵术语更长，但它包含了自学、复现和引用时真正需要的信息。
把本章用于自学时，可以按三遍阅读。第一遍只追踪对象和定义，确认每个术语到底指向文本、token、张量、样本、请求还是评估记录。第二遍追踪公式和伪代码，确认每一步改变了哪些状态。第三遍追踪实验和 caveat，确认哪些结论只在特定规模、数据或硬件条件下成立。
把本章用于复习时，则应反过来操作：先从一个失败案例或异常指标出发，再回到相应机制。loss 曲线异常可能指向数据、优化或模型容量；吞吐异常可能指向 kernel、通信或调度；benchmark 异常可能指向 contamination、prompt protocol 或采样设置。能够从异常反推机制，是掌握本章的实用标准。
\section{例题：把概念变成可计算检查}
\subsection{Serving latency 的分解}
一个用户请求从进入队列到 token 流式返回。
\begin{itemize}
\item 排队时间取决于调度和 batching。
\item prefill 处理 prompt，decode 循环逐 token 运行。
\item KV cache 管理、kernel 效率和输出长度共同影响 tail latency。
\end{itemize}
\textbf{结论。} serving 研究把模型结构、GPU kernels 和系统调度连成一条端到端路径。
这个例题的目的不是替代完整工程实现，而是给出一种最小推理形式：把问题转成变量，写出近似公式或伪代码，再说明近似何时失效。
\section{实验设计：从小例子到可复现结论}
学习 Guest Lecture: Dan Fu 与推理引擎研究 不能停在概念层。一个教材化结论至少需要四件事：可构造的输入、可观测的中间量、可比较的指标，以及清楚的失败条件。下面的实验设计不是要求训练大模型，而是要求读者用小规模程序复现本章机制。
\subsection{最小输入}
先构造只触发本章核心机制的小输入。例如 tokenizer 章节可以用几个包含 rare words 和 Unicode 字符的字符串；GPU 章节可以用一个 elementwise kernel 和一个 GEMM；data 章节可以用十几篇近重复文档。小输入的价值在于它让错误可见。
\subsection{资源账本}
对同一输入写下 FLOPs、bytes moved、显存峰值、通信量或样本数量的估算。估算不需要一开始就精确，但必须说明单位和近似假设。随后用 profiler、benchmark、validation loss 或人工检查去验证估算是否同量级。
\subsection{单变量消融}
一次只改变一个因素，例如 vocabulary size、head 数、batch size、filter threshold、reward scale、decoding temperature 或 GPU 数量。若多个变量同时变化，实验只能说明系统变了，不能说明机制是什么。
\subsection{失败条件}
最后要刻意触发失败：极端 context length、错误 dtype、过小 batch、污染 benchmark、低质量数据或不可靠 verifier。失败条件常常比成功曲线更能说明一个方法的边界。
\subsection{记录与报告}
实验记录应能让另一个读者复现同一结论。最少要写清楚输入构造、代码版本、随机种子、硬件、batch shape、数据版本、评价指标和异常处理方式。报告结论时不要只写“更快”或“更好”，而要说明快在哪里、好在哪个指标上，以及是否牺牲了内存、稳定性、数据质量或可解释性。
常见报告错误是把局部现象写成普遍规律。例如一次小模型实验里的 loss 改善，不等于更大模型或不同数据分布上也会改善；一次 kernel benchmark 的吞吐提升，不等于端到端 serving latency 一定降低。教材中的实验报告应始终把结论限定在可复现条件内。
\subsection{从小实验到大结论}
小实验的作用不是替代课程中的完整训练或系统实现，而是建立一条可检查的推理链。先在小输入上确认变量方向，再在中等规模上确认数量级，最后才讨论是否能外推到更大模型、更多 token 或更复杂 workload。缺少中间层级时，大结论通常只是在复述直觉。
因此，实验报告最好同时包含正例和反例：正例说明机制在受控条件下怎样工作，反例说明条件稍变时哪里失效。对 CS336 这样的课程，反例尤其重要，因为 language model 的错误常来自多个层面叠加：数据分布、优化稳定性、硬件瓶颈、评估协议和服务负载都可能同时改变观察结果。
还要记录资料版本与复现边界。公开视频、课程页、配套代码和 PDF 可能在不同时间更新；复现实验应写明使用的材料版本，并说明哪些结论来自课程事实，哪些是为了帮助理解而加入的解释性推导。这样才能把学习笔记变成可检查的技术记录。
若复现实验与课程结论不一致，首先检查输入、版本和评估协议，而不是立即修改模型假设。
\section{课程资料与回看路径}
教材正文已经把主要概念、公式和实现逻辑组织成连续叙述；本节只提供回到课程材料的阅读路径。读者复习时可以先完成前文练习，再按下面的时间段或资料组核对细节。
\subsection{视频回看路径}
\begin{longtable}{p{0.20\linewidth}p{0.46\linewidth}p{0.26\linewidth}}
\toprule
\textbf{时间} & \textbf{关键词} & \textbf{复习任务}\\
\midrule
00:00:05-00:06:23 & think, inference, horses, pretty, understand, code, around & 回看定义、例子、推导或 caveat 在该段中的位置。\\
00:17:45-00:23:43 & different, very, then, compute, prefill, decode, start & 回看定义、例子、推导或 caveat 在该段中的位置。\\
00:35:36-00:41:54 & your, attention, kernel, start, very, then, operation & 回看定义、例子、推导或 caveat 在该段中的位置。\\
00:47:41-00:53:39 & going, matrix, there's, look, then, it's, said & 回看定义、例子、推导或 caveat 在该段中的位置。\\
01:05:14-01:11:14 & think, then, your, size, cache, going, choices & 回看定义、例子、推导或 caveat 在该段中的位置。\\
\bottomrule
\end{longtable}
这些时间段用于定位视频中的讲解顺序。严肃引用时，应回到视频片段确认讲者表述、板书或幻灯片内容。
回看视频时，不必逐字抄录字幕。更有效的方法是把每个时间段判定为定义、例子、推导、代码解释或 caveat，再把它放回本章对应小节。
\subsection{课程资料阅读路径}
\begin{longtable}{p{0.14\linewidth}p{0.48\linewidth}p{0.30\linewidth}}
\toprule
\textbf{资料组} & \textbf{标题/页块} & \textbf{关键词}\\
\midrule
official\_01 & No official slide PDF or executable lecture script is linked for this public playlist ... & official, slide, executable, script, linked, public, playlist\\
\bottomrule
\end{longtable}
课程资料通常给出更稳定的标题、公式、图或代码结构；视频则补充动机和口头解释。两者合起来构成本章的事实边界。
阅读资料时，可以把每个资料组拆成三个对象：一个定义，一个公式或伪代码动作，一个边界条件。这样比逐字翻译页面或脚本更容易形成可用知识。
\section{实践说明、历史位置与风险}
Guest Lecture: Dan Fu 与推理引擎研究 在 CS336 的课程结构中连接基础机制和规模化问题。基础机制解释方法为什么成立；规模化问题检验它在更大模型、更长上下文、更复杂数据或真实 serving workload 下是否仍成立。
历史位置不等于模型年表。更重要的是识别问题如何迁移：早期方法解决小规模建模问题，现代方法常常解决同一问题在大规模数据、GPU 集群、长上下文或人类偏好约束下的新形态。
\begin{itemize}
\item 资料版本限制（version risk）：课程页、公开视频和配套资料可能在学期中更新；引用时应注明访问日期或资料版本。
\item 测量口径限制（measurement risk）：FLOPs、tokens/s、benchmark score、reward 等数字只有在协议完整时可比较。
\item 规模外推限制（scaling risk）：小模型或小 batch 的机制不必然外推到 frontier-scale。
\item 数据分布限制（data risk）：数据来源、过滤和去重策略会改变模型行为，也会改变 evaluation contamination 风险。
\item 系统风险（systems risk）：性能剖析结果依赖硬件、driver、kernel、shape 和 batch distribution；脱离 workload 的性能数字不可直接复用。
\end{itemize}
\section{复习要点}
下面的问题把本章内容压缩为可反复检查的条目。每个条目都应能回到前文的解释、公式、伪代码或例题。
\begin{longtable}{p{0.16\linewidth}p{0.38\linewidth}p{0.38\linewidth}}
\toprule
\textbf{条目} & \textbf{对象} & \textbf{检查动作}\\
\midrule
条目 1 & serving system（服务系统） & 定义它；写出一个最小例子；说明一个 failure mode；指出它影响哪个资源或评估账本。\\
条目 2 & inference engine（推理引擎） & 定义它；写出一个最小例子；说明一个 failure mode；指出它影响哪个资源或评估账本。\\
条目 3 & GPU kernel & 定义它；写出一个最小例子；说明一个 failure mode；指出它影响哪个资源或评估账本。\\
条目 4 & ThunderKittens & 定义它；写出一个最小例子；说明一个 failure mode；指出它影响哪个资源或评估账本。\\
条目 5 & decode loop（解码循环） & 定义它；写出一个最小例子；说明一个 failure mode；指出它影响哪个资源或评估账本。\\
条目 6 & state-space model (SSM) & 定义它；写出一个最小例子；说明一个 failure mode；指出它影响哪个资源或评估账本。\\
\bottomrule
\end{longtable}
\subsection{综合自测}
\begin{itemize}
\item 我能否不用课程原句，独立解释本章中心问题？
\item 我能否把本章至少一个公式改写成可运行的估算函数？
\item 我能否指出一个看似改进但实际转移成本的方法？
\item 我能否回到视频和课程资料，找到支持本章关键论断的位置？
\item 我能否设计一个最小 ablation，验证本章一个 caveat？
\end{itemize}
\section{总结与延伸}
\subsection{本章小结}
\begin{itemize}
\item Guest Lecture: Dan Fu 与推理引擎研究 的核心不是记忆术语，而是理解对象、变量和机制之间的关系。
\item 视频给出讲解顺序，课程资料提供可核对的公式、图、代码或页面结构。
\item 复习时应把每个术语连接到至少一个变量、一个实现动作和一个失败模式。
\end{itemize}
\subsection{延伸解释}
本章中的延伸解释用于连接课程材料中的概念，例如说明公式里的 proxy、解释系统 tradeoff，或把多个材料片段串成一个完整推理链。若需要严肃引用，应回到课程视频和配套资料核对原始表述。
\subsection{练习}
\begin{itemize}
\item 定义题：用中英双语解释本章任意五个重要术语，并说明它们属于 compute、memory、data、optimization、evaluation 还是 deployment 账本。
\item 推导题：选择本章一个公式，逐个解释符号、单位和近似假设，并说明哪个变量最难在真实系统中测量。
\item 实现题：把本章伪代码改写成一个最小 Python sanity check，不要求训练大模型，但要输出可检查的中间量。
\item 资料题：从“课程资料与回看路径”中选择一个视频时间段，回到原始视频核对讲者例子，并记录是否存在字幕误差。
\item 评估题：为本章方法设计一个 ablation 或 benchmark，说明控制变量、数据版本、指标和 failure criteria。
\item 边界题：说明本章一个结论在更长 context、更大 batch、更多 GPU 或更复杂数据分布下可能如何失效。
\item 综合题：把本章内容和前一章或后一章连接起来，写出一个端到端训练/推理 pipeline 中的依赖关系。
\item 批判题：指出一个容易被营销材料夸大的说法，并用本章的资源账本或评估协议反驳它。
\end{itemize}
\end{document}
